\documentclass[../../main.tex]{subfiles}

\graphicspath{{../../image/}{../../figure/}}

\begin{document}

\section{习题}\label{section:ch04-习题}

\begin{exercise}\label{exercise:ch04-01}
（★）考虑式\eqref{align:ch01-平方和误差}给出的平方和误差函数，其中函数$y(x,\boldsymbol{w})$由多项式\eqref{align:ch01-多项式模型}给出。证明能使该误差函数最小的系数$\boldsymbol{w}=\{w_i\}$是由下列线性方程组的解给出的：
\begin{align}
  \sum_{j=0}^{M}A_{ij}w_j=T_i
  \label{align:ch04-习题多项式方程}
\end{align}
其中
\begin{align}
  A_{ij}=\sum_{n=1}^{N}(x_n)^{i+j},\qquad T_i=\sum_{n=1}^{N}(x_n)^{i}t_n.
  \label{align:ch04-习题多项式系数}
\end{align}
\end{exercise}

\begin{exercise}\label{exercise:ch04-02}
（★）写出类似于式\eqref{align:ch04-习题多项式方程}的耦合线性方程组，其中系数$w_i$能使式\eqref{align:ch01-正则化误差}给出的正则化平方和误差函数最小化。
\end{exercise}

\begin{exercise}\label{exercise:ch04-03}
（★）证明定义如下的$\tanh$函数
\begin{align}
  \tanh(a)=\frac{\mathrm{e}^{a}-\mathrm{e}^{-a}}{\mathrm{e}^{a}+\mathrm{e}^{-a}}
  \label{align:ch04-习题tanh定义}
\end{align}
和式\eqref{align:ch04-sigmoid}定义的sigmoid函数有如下关系：
\begin{align}
  \tanh(a)=2\sigma(2a)-1
  \label{align:ch04-习题tanh关系}
\end{align}
并进一步证明sigmoid函数的一般线性组合形式
\begin{align}
  y(x,\boldsymbol{w})=w_0+\sum_{j=1}^{M}w_j\sigma\Bigl(\frac{x-\mu_j}{s}\Bigr)
  \label{align:ch04-习题sigmoid组合}
\end{align}
等价于如下形式的$\tanh$函数的线性组合：
\begin{align}
  y(x,\boldsymbol{u})=u_0+\sum_{j=1}^{M}u_j\tanh\Bigl(\frac{x-\mu_j}{2s}\Bigr)
  \label{align:ch04-习题tanh组合}
\end{align}
最后找到将新参数$\{u_1,\cdots,u_M\}$与原始参数$\{w_1,\cdots,w_M\}$关联起来的表达式。
\end{exercise}

\begin{exercise}\label{exercise:ch04-04}
（★★★）证明矩阵
\begin{align}
  \boldsymbol{\Phi}(\boldsymbol{\Phi}^{\mathsf{T}}\boldsymbol{\Phi})^{-1}\boldsymbol{\Phi}^{\mathsf{T}}
  \label{align:ch04-习题投影矩阵}
\end{align}
可将任何向量$\boldsymbol{v}$投影到$\boldsymbol{\Phi}$的列所张成的空间中。使用此结果证明，最小二乘解式\eqref{align:ch04-最小二乘解}对应向量$\boldsymbol{t}$在流形$\mathcal{S}$上的正交投影（见\cref{figure:ch04-最小二乘几何}）。
\end{exercise}

\begin{exercise}\label{exercise:ch04-05}
（★）考虑一个每个数据点$t_n$都与加权因子$r_n>0$相关联的数据集，平方和误差函数变为
\begin{align}
  E_D(\boldsymbol{w})=\frac{1}{2}\sum_{n=1}^{N}r_n\bigl\{t_n-\boldsymbol{w}^{\mathsf{T}}\boldsymbol{\phi}(\boldsymbol{x}_n)\bigr\}^2 .
  \label{align:ch04-习题加权误差}
\end{align}
给出最小化该误差函数的解$\boldsymbol{w}^*$的表达式。从与数据相关的噪声方差和复制的数据点两方面，给出加权平方和误差函数的两种替代解释。
\end{exercise}

\begin{exercise}\label{exercise:ch04-06}
（★）通过将式\eqref{align:ch04-正则化总误差}对$\boldsymbol{w}$的梯度设置为$\boldsymbol{0}$，证明在线性回归中，正则化平方和误差函数的精确最小值由式\eqref{align:ch04-正则化解}给出。
\end{exercise}

\begin{exercise}\label{exercise:ch04-07}
（★★）假设多元目标变量$\boldsymbol{t}$的线性基函数回归模型具有下式给出的Gauss分布：
\begin{align}
  p(\boldsymbol{t}\mid\boldsymbol{W},\boldsymbol{\Sigma})=\mathcal{N}\bigl(\boldsymbol{t}\mid\boldsymbol{y}(\boldsymbol{x},\boldsymbol{W}),\boldsymbol{\Sigma}\bigr)
  \label{align:ch04-习题多元Gauss}
\end{align}
其中
\begin{align}
  \boldsymbol{y}(\boldsymbol{x},\boldsymbol{W})=\boldsymbol{W}^{\mathsf{T}}\boldsymbol{\phi}(\boldsymbol{x}) .
  \label{align:ch04-习题多元均值}
\end{align}
训练数据集由输入基向量$\boldsymbol{\phi}(\boldsymbol{x}_n)$和对应的目标向量$\boldsymbol{t}_n$组成，其中$n=1,\cdots,N$。证明参数矩阵$\boldsymbol{W}$的最大似然解$\boldsymbol{W}_{\mathrm{ML}}$具有如下性质：每列由形如式\eqref{align:ch04-最小二乘解}的表达式给出，其是各向同性噪声分布的解。注意，它独立于协方差矩阵$\boldsymbol{\Sigma}$。证明$\boldsymbol{\Sigma}$的最大似然解由下式给出：
\begin{align}
  \boldsymbol{\Sigma}=\frac{1}{N}\sum_{n=1}^{N}\bigl(\boldsymbol{t}_n-\boldsymbol{W}_{\mathrm{ML}}^{\mathsf{T}}\boldsymbol{\phi}(\boldsymbol{x}_n)\bigr)\bigl(\boldsymbol{t}_n-\boldsymbol{W}_{\mathrm{ML}}^{\mathsf{T}}\boldsymbol{\phi}(\boldsymbol{x}_n)\bigr)^{\mathsf{T}} .
  \label{align:ch04-习题协方差ML}
\end{align}
\end{exercise}

\begin{exercise}\label{exercise:ch04-08}
（★）考虑将单个目标变量$t$的平方损失函数式\eqref{align:ch04-平方损失期望}推广到下式给出的多目标变量（用向量$\boldsymbol{t}$表示），则有
\begin{align}
  \mathbb{E}\bigl[L(\boldsymbol{t},\boldsymbol{f}(\boldsymbol{x}))\bigr]=\iint\bigl\|\boldsymbol{f}(\boldsymbol{x})-\boldsymbol{t}\bigr\|^2p(\boldsymbol{x},\boldsymbol{t})\,\mathrm{d}\boldsymbol{x}\,\mathrm{d}\boldsymbol{t} .
  \label{align:ch04-习题多元损失}
\end{align}
使用变分演算，证明使预期损失最小化的函数$\boldsymbol{f}(\boldsymbol{x})$由式\eqref{align:ch04-习题多元最优}给出：
\begin{align}
  \boldsymbol{f}(\boldsymbol{x})=\mathbb{E}_{\boldsymbol{t}}[\boldsymbol{t}\mid\boldsymbol{x}] .
  \label{align:ch04-习题多元最优}
\end{align}
\end{exercise}

\begin{exercise}\label{exercise:ch04-09}
（★）通过对式\eqref{align:ch04-习题多元损失}中的平方项进行展开，推导类似于式\eqref{align:ch04-平方损失分解}的结果。进一步地，证明使多元目标变量$\boldsymbol{t}$的预期平方损失最小化的函数$\boldsymbol{f}(\boldsymbol{x})$同样以式\eqref{align:ch04-习题多元最优}的形式由$\boldsymbol{t}$的条件期望给出。
\end{exercise}

\begin{exercise}\label{exercise:ch04-10}
（★★）类比式\eqref{align:ch04-平方损失分解}，通过扩展式\eqref{align:ch04-习题多元损失}来重新推导结果式\eqref{align:ch04-习题多元最优}。
\end{exercise}

\begin{exercise}\label{exercise:ch04-11}
（★★）以下分布
\begin{align}
  p(x\mid\sigma^2,q)=\frac{q}{2(2\sigma^2)^{1/q}\Gamma(1/q)}\exp\Bigl(-\frac{|x|^q}{2\sigma^2}\Bigr)
  \label{align:ch04-习题广义Gauss}
\end{align}
是一元Gauss分布的推广形式。其中$\Gamma(x)$是由下式定义的伽马函数：
\begin{align}
  \Gamma(x)=\int_{-\infty}^{\infty}u^{x-1}\mathrm{e}^{-u}\,\mathrm{d}u .
  \label{align:ch04-习题伽马函数}
\end{align}
证明此分布已经归一化，从而有
\begin{align}
  \int_{-\infty}^{\infty}p(x\mid\sigma^2,q)\,\mathrm{d}x=1,
  \label{align:ch04-习题归一化}
\end{align}
且当$q=2$时，它退化为Gauss分布。考虑一个目标变量由$t=y(\boldsymbol{x},\boldsymbol{w})+\varepsilon$给出的回归模型，$\varepsilon$是从分布\eqref{align:ch04-习题广义Gauss}中产生的随机噪声变量。证明对于输入向量$X=\{\boldsymbol{x}_1,\cdots,\boldsymbol{x}_N\}$和相应的目标变量$\boldsymbol{t}=(t_1,\cdots,t_N)^{\mathsf{T}}$的数据集，定义在$\boldsymbol{w}$和$\sigma^2$上的对数似然函数由下式给出：
\begin{align}
  \ln p(\boldsymbol{t}\mid X,\boldsymbol{w},\sigma^2)=-\frac{1}{2\sigma^2}\sum_{n=1}^{N}\bigl|y(\boldsymbol{x}_n,\boldsymbol{w})-t_n\bigr|^q-\frac{N}{q}\ln(2\sigma^2)+\text{const},
  \label{align:ch04-习题广义对数似然}
\end{align}
其中“const”表示独立于$\boldsymbol{w}$和$\sigma^2$的常数项。注意，作为$\boldsymbol{w}$的函数，该对数似然函数是\cref{section:ch04-决策理论}中的$L_q$误差函数。
\end{exercise}

\begin{remark}（重构新增：原文问题说明）
原页把伽马函数印成$\Gamma(x)=\int_{-\infty}^{\infty}u^{x-1}\mathrm{e}^{-u}\,\mathrm{d}u$（下限为$-\infty$），这一写法有误：当$u\to-\infty$且$x-1$非整数时$u^{x-1}$无定义，该积分并不收敛。本译本在正文中照录了这一印刷形式，但在归一化证明中采用了下限为$0$的标准定义，即$\Gamma(x)=\int_{0}^{\infty}u^{x-1}\mathrm{e}^{-u}\,\mathrm{d}u$；归一化证明\eqref{align:ch04-习题归一化}正是在这一形式下进行的。除此之外原题的条件与结论不受影响。
\end{remark}

\begin{exercise}\label{exercise:ch04-12}
（★★）考虑由$L_q$损失函数定义的回归问题的期望损失式\eqref{align:ch04-Minkowski损失}。写出在最小化$\mathbb{E}[L_q]$时，$y(\boldsymbol{x})$必须满足的条件。证明当$q=1$时，该解表示条件中位数，即函数$y(\boldsymbol{x})$满足$t<y(\boldsymbol{x})$的概率与$t\geqslant y(\boldsymbol{x})$的概率相同；并证明当$q\to0$时最小期望损失$L_q$由条件众数给出，即函数$y(\boldsymbol{x})$等于对每个$\boldsymbol{x}$而言，能使$p(t\mid\boldsymbol{x})$最大化的$t$值。
\end{exercise}

\end{document}
